\documentclass[12pt,conference,onecolumn]{IEEEtran}

\usepackage[hidelinks]{hyperref}

\title{Title here: Internship at Commvault}
\author{%
\IEEEauthorblockN{Ethan Fuks}\IEEEauthorblockA{Science \& Engineering\\Manalapan High School\\Englishtown, NJ\\\href{mailto:426efuks@frhsd.com}{426efuks@frhsd.com}}}
\date{June 16, 2026}

\newcommand{\keywords}{Commvault, malware, MISP, hash function, MalwareBazaar, abuse.ch, internship}

\usepackage{hyperref}
\makeatletter
\AtBeginDocument{
\hypersetup{%
pdftitle={\@title},
pdfauthor={Ethan Fuks},
pdfkeywords={\keywords}}}
\makeatother

\begin{document}
\maketitle 

\begin{abstract}
During my internship at Commvault, I have created a tool that recurrently checks for malware in a user’s data by using MISP to obtain newly discovered malware hashes and comparing these against the user. A hash is a unique way of representing a piece of data by encoding it where small changes in the input create big unpredictable changes in the output. This creates a unique ``name'' for each file that can be used to share found malware and compare against it. MISP is a software that organizes and collects different feeds of data from various sources including \href{https://bazaar.abuse.ch}{MalwareBazaar} and \href{https://abuse.ch}{abuse.ch}, among many others. I created a dummy system to test with 100 files of randomized text and a list of hashes of these was created. Some were selected as malware and were checked against the list. I set up a local instance of MISP on a Commvault virtual machine that is perpetually running and able to be queried for found malware. This is then integrated into the dummy system to recurrently check for malware. I created a write-up for my process of setting up and using MISP for Commvault to use in the future. 
\end{abstract}

\begin{IEEEkeywords}
\keywords
\end{IEEEkeywords}

\end{document}
